\section{Kesimpulan}

Nadina menyediakan alur penandatanganan dan verifikasi \textit{release} untuk tiga profil
algoritma, yaitu RSASSA-PSS dengan RSA 3072 bit, ECDSA P-256, dan Ed25519. Program dibangun
hanya dari pustaka OTP, dan didistribusikan sebagai image \textit{container} multi-arch yang dikunci
per \textit{digest}. Di atasnya terdapat \textit{trust store} yang dikelola pengguna, verifikasi
tujuh langkah dengan hasil JSON dan exit code yang tegas, serta pengikatan SBOM ke artefak rilis.

Ada 250 tes pada mesin pengujian, termasuk skenario U01 sampai U14 untuk 
ketiga profil, \textit{property test}, dan interoperabilitas dengan
OpenSSL untuk ketiga profil. Pada pipeline pemakai, skenario C01 sampai C05 lulus pada image
rilis \texttt{0.1.0}. Skenario negatif C02 sampai C04 menghasilkan exit code yang sesuai dan tidak
membuat penanda deploy, sehingga penolakan benar-benar menghentikan alur. Benchmark pada satu
mesin menunjukkan bahwa \textit{sign} RSA-PSS sekitar 68 kali lebih lambat daripada ECDSA,
sedangkan biaya proses pada alur ujung ke ujung didominasi oleh boot BEAM.

Wrapper SBOM menyimpan byte asli dalam Base64 agar
\textit{signature} mencakup byte yang tepat dan keluaran \textit{extract} identik dengan keluaran
kakas asalnya. Serialisasi \textit{canonical} manifest dibandingkan byte demi byte saat verifikasi,
sehingga byte yang ditandatangani pasti sama dengan byte yang diperiksa. Exit code dijaga agar
tidak pernah berubah menjadi 0 di antara proses dan pipeline.

Rancangan  yang kami buat tidak mencegah \textit{rollback attack} tanpa \textit{trusted state}
tambahan. Tidak ada \textit{transparency log}, dan tanda tangan hanya memeriksa asal dan keutuhan,
bukan keamanan perangkat lunak. Isi SBOM tidak diperiksa kelengkapannya. Pengujian interoperabilitas
di dalam image Alpine dilewati karena image tidak memuat CLI \texttt{openssl}, sehingga tes tersebut
dijalankan di mesin host. Benchmark hanya dilakukan pada satu mesin.
